\chapter{向量积分}

在本章中，若 F 表示（R 或 C 上的）Hausdorff 局部凸向量空间，我们用 \(F'\) 表示它的对偶，用 \(F''\) 表示它的双对偶，用 \(F'^*\) 表示 \(F'\) 的代数对偶（即 \(F'\) 上全体线性形式所成的空间）；\(F''\) 是 \(F'^*\) 的线性子空间，而 F（作为不带拓扑的向量空间）可以与 \(F''\) 的一个线性子空间等同。我们用 \(F_\sigma\) 表示赋以弱化拓扑 \(\sigma(F, F')\) 的向量空间 F；“弱的”、“弱地”等定语即指这一拓扑。

在本章中，T 表示局部紧空间，\(\mathcal{K}_{\mathbf{R}}(\mathrm{T})\) 或 \(\mathcal{K}(\mathrm{T})\)（相应地 \(\mathcal{K}_{\mathbf{C}}(\mathrm{T})\)）表示 T 上连续且具有紧支集的实（相应地复）值函数所成的向量空间；对 T 的每个子集 A，\(\mathcal{K}(\mathrm{T},\mathrm{A})\)（相应地 \(\mathcal{K}_{\mathbf{C}}(\mathrm{T},\mathrm{A})\)）表示 \(\mathcal{K}(\mathrm{T})\)（相应地 \(\mathcal{K}_{\mathbf{C}}(\mathrm{T})\)）中由支集含于 A 的函数所构成的子空间。除非另有明确说明，空间 \(\mathcal{K}(\mathrm{T})\)（相应地 \(\mathcal{K}_{\mathbf{C}}(\mathrm{T})\)）都赋以在各子空间 \(\mathcal{K}(\mathrm{T},\mathrm{K})\)（相应地 \(\mathcal{K}_{\mathbf{C}}(\mathrm{T},\mathrm{K})\)）上一致收敛拓扑的直接极限拓扑，其中 K 跑遍 T 的紧子集全体。

我们提醒，这个拓扑比一致收敛拓扑更细，因而是 Hausdorff 的；它在每个 \(\mathcal{K}(\mathrm{T},\mathrm{K})\)（相应地 \(\mathcal{K}_{\mathbf{C}}(\mathrm{T},\mathrm{K})\)）上诱导一致收敛拓扑（TVS, II, §4, No.~4, 注）。空间 \(\mathcal{K}_{\mathbf{C}}(\mathrm{T})\) 可以与由 \(\mathcal{K}(\mathrm{T})\) 经把标量域从 R 扩张到 C 所得的空间等同。说 \(\mathcal{K}(\mathrm{T})\) 上的一个线性形式是测度，按定义即指它是连续的（Ch. III, §1, No.~3, 定义 2）。

